\section{Worked Examples}

\subsection{Example MicroSim: A Bouncing Ball for Newtonian Mechanics}

Consider a MicroSim that visualizes a bouncing ball under gravity, targeting
the concept \textit{Newton's Second Law} within a physics textbook. The base
variant exposes a single control: initial drop height. Several evolutionary
hypotheses can be generated and tested:

\begin{itemize}
    \item \textbf{Variant A (graph overlay):} Add a real-time plot of velocity
          versus time alongside the animation. \textit{Hypothesis: connecting
          the visual motion to its quantitative representation improves transfer
          to symbolic problem-solving.}
    \item \textbf{Variant B (parameter reduction):} For novice learners, hide
          all controls except one, reducing extraneous cognitive load.
          \textit{Hypothesis: fewer simultaneous variables improves initial
          conceptual grasp.}
    \item \textbf{Variant C (prediction prompt):} Before running the animation,
          ask the learner to predict where the ball will be after one second.
          \textit{Hypothesis: active prediction (a generative learning
          strategy) improves retention.}
\end{itemize}

Each variant is deployed, its composite fitness measured, and the winner
promoted. Note that Variant C tests a pedagogical strategy (prediction-based
active learning) that generalizes across many MicroSims; a win here might
propagate as a mutation operator applied to other concepts, mirroring how the
DGM discovered general-purpose improvements such as solution-ranking that
transferred across tasks.

\subsection{Example MicroSim: Compound Interest for Personal Finance}

A compound interest visualizer in a personal finance textbook lets learners
adjust principal, interest rate, and time horizon, and observe the resulting
growth curve. This MicroSim illustrates an important archival principle: a
variant that performs poorly for one audience may be the ideal stepping stone
for another. A visually dense variant showing multiple overlaid scenarios might
overwhelm high-school learners but excel for adult learners with financial
background. The archive retains both, and selection becomes
audience-conditioned---a natural extension of the concept-graph search space
into a learner-segment dimension.

\subsection{Example Textbook Domains}

The framework applies across the disciplinary breadth typical of an
intelligent-textbook portfolio. MicroSim-rich quantitative domains such as
physics and biology at the secondary and International Baccalaureate level are
especially strong candidates, because the abundance of simulatable phenomena
gives the evolutionary process a large space of meaningful mutations. Computing
curricula such as algorithms and data structures are equally promising: an
animated sorting or graph-traversal MicroSim has many tunable pedagogical
dimensions (animation speed, step annotation, comparison highlighting) that
map naturally onto evolutionary parameters. History and other narrative domains
benefit less from simulation but can still evolve their explanatory text and
concept-graph structure.

\subsection{Skill-Level Self-Improvement}

A subtle but powerful application targets not the MicroSims themselves but the
\textit{content-generation skills} that produce them. Following the DGM's
self-referential principle---that improving the codebase also improves the
ability to improve the codebase---the prompts and templates used to generate
MicroSims can themselves be treated as an evolving population. Variant prompts
are scored by the average fitness of the MicroSims they produce, and
high-performing prompts are retained and mutated. Over time, the generator
improves, a genuine second-order effect that compounds the gains from
content-level evolution.
